\chapter*{References}
\addcontentsline{toc}{chapter}{References}

Athey, S., \& Imbens, G. (2016). Recursive partitioning for heterogeneous causal effects. Proceedings of the National Academy of Sciences, 113(27), 7353-7360. \href{https://www.pnas.org/doi/full/10.1073/pnas.1510489113}{https://doi.org/10.1073/pnas.1510489113}

Bray, F., Laversanne, M., Sung, H., Ferlay, J., Siegel, R. L., Soerjomataram, I., \& Jemal, A. (2024). Global cancer statistics 2022: GLOBOCAN estimates of incidence and mortality worldwide for 36 cancers in 185 countries. CA: A Cancer Journal for Clinicians, 74(3), 229-263. \href{https://acsjournals.onlinelibrary.wiley.com/doi/10.3322/caac.21834}{https://doi.org/10.3322/caac.21834}

Chen, T., \& Guestrin, C. (2016). XGBoost: A scalable tree boosting system. In Proceedings of the 22nd ACM SIGKDD International Conference on Knowledge Discovery and Data Mining (pp. 785-794). \href{https://dl.acm.org/doi/10.1145/2939672.2939785}{https://doi.org/10.1145/2939672.2939785}

Chung, P., Swaminathan, A., et al. (2025). VeriFact: Verifying facts in LLM-generated clinical text with electronic health records. arXiv:2501.16672. \href{https://arxiv.org/abs/2501.16672}{https://arxiv.org/abs/2501.16672}

Colombo, N., Sessa, C., du Bois, A., Ledermann, J., McCluggage, W. G., McNeish, I., \& Vergote, I. (2019). ESMO-ESGO consensus conference recommendations on ovarian cancer: Pathology and molecular biology, early and advanced stages, borderline tumours and recurrent disease. Annals of Oncology, 30(5), 672-705. \href{https://www.annalsofoncology.org/article/S0923-7534(19)31162-7/fulltext}{https://doi.org/10.1093/annonc/mdz062}

Cox, D. R. (1972). Regression models and life-tables. Journal of the Royal Statistical Society, Series B (Methodological), 34(2), 187-220.

Curth, A., Svensson, D., Weatherall, J., \& van der Schaar, M. (2021a). Really doing great at estimating CATE? A critical look at ML benchmarking practices in treatment effect estimation. In Proceedings of the 35th Conference on Neural Information Processing Systems Datasets and Benchmarks Track.

Curth, A., Lee, C., \& van der Schaar, M. (2021b). SurvITE: Learning heterogeneous treatment effects from time-to-event data. In Advances in Neural Information Processing Systems 34 (pp. 26740-26753).

Dud\'{i}k, M., Langford, J., \& Li, L. (2011). Doubly robust policy evaluation and learning. In Proceedings of the 28th International Conference on Machine Learning (pp. 1097-1104).

European Parliament and Council. (2024). Regulation (EU) 2024/1689 of the European Parliament and of the Council of 13 June 2024 laying down harmonised rules on artificial intelligence (Artificial Intelligence Act). Official Journal of the European Union, L 2024/1689. \url{https://eur-lex.europa.eu/eli/reg/2024/1689/oj}

Gao, H., Korn, J. M., Ferretti, S., Monahan, J. E., Wang, Y., Singh, M., Zhang, C., Schnell, C., Yang, G., Zhang, Y., Balbin, O. A., Barbe, S., Cai, H., Casey, F., Chatterjee, S., Chiang, D. Y., Chuai, S., Cogan, S. M., Collins, S. D., ... Sellers, W. R. (2015). High-throughput screening using patient-derived tumor xenografts to predict clinical trial drug response. Nature Medicine, 21(11), 1318-1325. \href{https://www.nature.com/articles/nm.3954}{https://doi.org/10.1038/nm.3954}

Goddard, K., Roudsari, A., \& Wyatt, J. C. (2012). Automation bias: A systematic review of frequency, effect mediators, and mitigators. Journal of the American Medical Informatics Association, 19(1), 121-127. \href{https://academic.oup.com/jamia/article-abstract/19/1/121/732254?redirectedFrom=fulltext}{https://doi.org/10.1136/amiajnl-2011-000089}

Hanker, L. C., Loibl, S., Burchardi, N., Pfisterer, J., Meier, W., Pujade-Lauraine, E., Ray-Coquard, I., Sehouli, J., Harter, P., \& du Bois, A. (2012). The impact of second to sixth line therapy on survival of relapsed ovarian cancer after primary taxane/platinum-based therapy. Annals of Oncology, 23(10), 2605-2612. \href{https://www.annalsofoncology.org/article/S0923-7534(19)37977-3/fulltext}{https://doi.org/10.1093/annonc/mds203}

Holland, P. W. (1986). Statistics and causal inference. Journal of the American Statistical Association, 81(396), 945-960. \href{https://www.tandfonline.com/doi/abs/10.1080/01621459.1986.10478354}{https://doi.org/10.1080/01621459.1986.10478354}

Imbens, G. W., \& Rubin, D. B. (2015). Causal inference for statistics, social, and biomedical sciences: An introduction. Cambridge University Press.

Ishwaran, H., Kogalur, U. B., Blackstone, E. H., \& Lauer, M. S. (2008). Random survival forests. Annals of Applied Statistics, 2(3), 841-860. \href{https://projecteuclid.org/journals/annals-of-applied-statistics/volume-2/issue-3/Random-survival-forests/10.1214/08-AOAS169.full}{https://doi.org/10.1214/08-AOAS169}

Kennedy, E. H. (2023). Towards optimal doubly robust estimation of heterogeneous causal effects. Electronic Journal of Statistics, 17(2), 3008-3049. \href{https://doi.org/10.1214/23-EJS2157}{https://doi.org/10.1214/23-EJS2151}

Kilim, O., Olar, A., Biricz, A., Madaras, L., Pollner, P., Szallasi, Z., Sztupinszki, Z., \& Csabai, I. (2025). Histopathology and proteomics are synergistic for high-grade serous ovarian cancer platinum response prediction. npj Precision Oncology, 9, 27. \href{https://www.nature.com/articles/s41698-025-00808-w}{https://doi.org/10.1038/s41698-025-00808-w}

Kunzel, S. R., Sekhon, J. S., Bickel, P. J., \& Yu, B. (2019). Metalearners for estimating heterogeneous treatment effects using machine learning. Proceedings of the National Academy of Sciences, 116(10), 4156-4165. \href{https://www.pnas.org/doi/full/10.1073/pnas.1804597116}{https://doi.org/10.1073/pnas.1804597116}

Landry, L. G., Ali, N., Williams, D. R., Rehm, H. L., \& Bonham, V. L. (2018). Lack of diversity in genomic databases is a barrier to translating precision medicine research into practice. Health Affairs, 37(5), 780-785. \href{https://www.healthaffairs.org/doi/10.1377/hlthaff.2017.1595}{https://doi.org/10.1377/hlthaff.2017.1595}

Louizos, C., Shalit, U., Mooij, J., Sontag, D., Zemel, R., \& Welling, M. (2017). Causal effect inference with deep latent-variable models. In Advances in Neural Information Processing Systems 30 (pp. 6446-6456).

Lundberg, S. M., \& Lee, S.-I. (2017). A unified approach to interpreting model predictions. In Advances in Neural Information Processing Systems 30 (pp. 4765-4774).

Lundberg, S. M., Erion, G., Chen, H., DeGrave, A., Prutkin, J. M., Nair, B., Katz, R., Himmelfarb, J., Bansal, N., \& Lee, S.-I. (2020). From local explanations to global understanding with explainable AI for trees. Nature Machine Intelligence, 2(1), 56-67. \href{https://www.nature.com/articles/s42256-019-0138-9}{https://doi.org/10.1038/s42256-019-0138-9}

Nie, X., \& Wager, S. (2021). Quasi-oracle estimation of heterogeneous treatment effects. Biometrika, 108(2), 299-319. \href{https://academic.oup.com/biomet/article-abstract/108/2/299/5911092?redirectedFrom=fulltext}{https://doi.org/10.1093/biomet/asaa076}

Obermeyer, Z., \& Emanuel, E. J. (2016). Predicting the future - big data, machine learning, and clinical medicine. New England Journal of Medicine, 375(13), 1216-1219. \href{https://www.nejm.org/doi/10.1056/NEJMp1606181}{https://doi.org/10.1056/NEJMp1606181}

Obermeyer, Z., Powers, B., Vogeli, C., \& Mullainathan, S. (2019). Dissecting racial bias in an algorithm used to manage the health of populations. Science, 366(6464), 447-453. \href{https://www.science.org/doi/10.1126/science.aax2342}{https://doi.org/10.1126/science.aax2342}

Pauli, C., Hopkins, B. D., Prandi, D., Shaw, R., Fedrizzi, T., Sboner, A., Sailer, V., Augello, M., Puca, L., Rosati, R., McNary, T. J., Churakova, Y., Cheung, C., Triscott, J., Pisapia, D., Rao, R., Mosquera, J. M., Robinson, B., Faltas, B. M., ... Rubin, M. A. (2017). Personalized in vitro and in vivo cancer models to guide precision medicine. Cancer Discovery, 7(5), 462-477. \href{https://aacrjournals.org/cancerdiscovery/article/7/5/462/6239/Personalized-In-Vitro-and-In-Vivo-Cancer-Models-to}{https://doi.org/10.1158/2159-8290.CD-16-1154}

Pearl, J. (2009). Causality: Models, reasoning, and inference (2nd ed.). Cambridge University Press.

Prigerson, H. G., Bao, Y., Shah, M. A., Paulk, M. E., LeBlanc, T. W., Schneider, B. J., Garrido, M. M., Reid, M. C., Berlin, D. A., Adelson, K. B., Neugut, A. I., \& Maciejewski, P. K. (2015). Chemotherapy use, performance status, and quality of life at the end of life. JAMA Oncology, 1(6), 778-784. \href{https://jamanetwork.com/journals/jamaoncology/fullarticle/2398177}{https://doi.org/10.1001/jamaoncol.2015.2378}

Reyes, M., Meier, R., Pereira, S., Silva, C. A., Dahlweid, F.-M., von Tengg-Kobligk, H., Summers, R. M., \& Wiest, R. (2020). On the interpretability of artificial intelligence in radiology: Challenges and opportunities. Radiology: Artificial Intelligence, 2(3), e190043. \href{https://pubs.rsna.org/doi/10.1148/ryai.2020190043}{https://doi.org/10.1148/ryai.2020190043}

Robinson, P. M. (1988). Root-N-consistent semiparametric regression. Econometrica, 56(4), 931-954. \href{https://www.jstor.org/stable/1912705?origin=crossref}{https://doi.org/10.2307/1912705}

Rubin, D. B. (1974). Estimating causal effects of treatments in randomized and nonrandomized studies. Journal of Educational Psychology, 66(5), 688-701. \href{https://psycnet.apa.org/doiLanding?doi=10.1037\%2Fh0037350}{https://doi.org/10.1037/h0037350}

Shalit, U., Johansson, F. D., \& Sontag, D. (2017). Estimating individual treatment effect: Generalization bounds and algorithms. In Proceedings of the 34th International Conference on Machine Learning (pp. 3076-3085).

Shenoy, N., Bhagat, T. D., Bhagat, A., \& Verma, A. (2020). Predicting ovarian cancer patients\textquotesingle{} clinical response to platinum-based chemotherapy by their tumor proteomic signatures. Journal of Proteome Research, 19(10), 4194-4206. \href{https://pubs.acs.org/doi/10.1021/acs.jproteome.0c00525}{https://doi.org/10.1021/acs.jproteome.0c00525}

Siegel, R. L., Giaquinto, A. N., \& Jemal, A. (2024). Cancer statistics, 2024. CA: A Cancer Journal for Clinicians, 74(1), 12-49. \href{https://acsjournals.onlinelibrary.wiley.com/doi/10.3322/caac.21820}{https://doi.org/10.3322/caac.21820}

Tayebi Arasteh, S., Han, T., Lotfinia, M., Kuhl, C., Kather, J. N., Truhn, D., \& Nebelung, S. (2024). Large language models streamline automated machine learning for clinical studies. Nature Communications, 15, 1603. \href{https://www.nature.com/articles/s41467-024-45879-8}{https://doi.org/10.1038/s41467-024-45879-8}

TCGA Research Network. (2011). Integrated genomic analyses of ovarian carcinoma. Nature, 474(7353), 609-615. \href{https://www.nature.com/articles/nature10166}{https://doi.org/10.1038/nature10166}

Thirunavukarasu, A. J., Ting, D. S. J., Elangovan, K., Gutierrez, L., Tan, T. F., \& Ting, D. S. W. (2023). Large language models in medicine. The Lancet Digital Health, 5(7), e384-e395. \href{https://www.nature.com/articles/s41591-023-02448-8}{https://doi.org/10.1016/S2589-7500(23)00080-5}

van Buuren, S., \& Groothuis-Oudshoorn, K. (2011). mice: Multivariate imputation by chained equations in R. Journal of Statistical Software, 45(3), 1-67. \href{https://www.jstatsoft.org/article/view/v045i03}{https://doi.org/10.18637/jss.v045.i03}

VanderWeele, T. J., \& Ding, P. (2017). Sensitivity analysis in observational research: Introducing the E-value. Annals of Internal Medicine, 167(4), 268-274. \href{https://doi.org/10.7326/M16-2607}{https://doi.org/10.7326/M16-2607}

Wager, S., \& Athey, S. (2018). Estimation and inference of heterogeneous treatment effects using random forests. Journal of the American Statistical Association, 113(523), 1228-1242. \href{https://www.tandfonline.com/doi/full/10.1080/01621459.2017.1319839}{https://doi.org/10.1080/01621459.2017.1319839}

Wohlin, C., Runeson, P., Host, M., Ohlsson, M. C., Regnell, B., \& Wesslen, A. (2012). Experimentation in software engineering. Springer.

Wright, A. A., Keating, N. L., Ayanian, J. Z., Chrischilles, E. A., Kahn, K. L., Ritchie, C. S., Weeks, J. C., Earle, C. C., \& Landrum, M. B. (2016). Family perspectives on aggressive cancer care near the end of life. JAMA, 315(3), 284-292. \href{https://jamanetwork.com/journals/jama/fullarticle/2482326}{https://doi.org/10.1001/jama.2015.18604}

Xiang, Q., Zhang, J., \& Feng, B. (2026). Using large language models for sensitivity analysis in causal inference: Case studies on Cornfield inequality and E-value. arXiv:2603.14273. \href{https://arxiv.org/abs/2603.14273}{https://arxiv.org/abs/2603.14273}

Yin, C., Yang, J., Wang, F., \& Sun, S. (2022). Prediction model of platinum response and prognosis in ovarian cancer based on multi-omics data. Frontiers in Oncology, 12. (Manual verification recommended at submission time.)
